\section{Анализ предметной области и требований к базе знаний}
\label{sec:analysis}

\subsection{Характеристика предметной области}
\label{subsec:domain-characteristics}

Предметная область курсового проекта — \textit{«Фундаментальные основы Computer Science»} — представляет собой систематизированное множество базовых концепций, принципов, алгоритмов и паттернов, составляющих теоретический и практический базис для успешного прохождения технических собеседований в сфере разработки программного обеспечения. Данная область является первым уровнем (\textit{Level 1}) в трёхуровневой архитектуре базы знаний интеллектуального помощника \textit{TechInterview} и служит фундаментом для двух других предметных областей: «Проектирование и архитектура систем» и «Процесс собеседования и soft skills»..

\subsubsection*{Типы единиц знаний предметной области}

Для представления знаний в базе знаний определены шесть основных типов единиц информации, каждый из которых выполняет специфическую роль в процессе подготовки к собеседованиям. Типы единиц знаний, их коды и описания приведены в таблице~\ref{tab:knowledge-units}.

\begin{table}[h]
\centering
\caption{Типы единиц знаний, специфичных для предметной области}
\label{tab:knowledge-units}
\begin{tabular}{|l|l|p{6cm}|}
\hline
\textbf{Тип} & \textbf{Код} & \textbf{Описание} \\ \hline
Карточка концепции & CONCEPT & Теоретическая концепция CS (структура данных, алгоритм, принцип) \\ \hline
Карточка задачи & PROBLEM & Алгоритмическая задача с решением и разбором \\ \hline
Паттерн решения & PATTERN & Типовой подход к классу задач \\ \hline
Карточка вопроса & QUESTION & Вопрос, задаваемый на собеседовании, с эталонным ответом \\ \hline
Шпаргалка & CHEATSHEET & Компактная сводка по теме для быстрого повторения \\ \hline
Типичная ошибка & PITFALL & Распространённая ошибка кандидата с объяснением \\ \hline
\end{tabular}
\end{table}

Каждый тип едини знаний характеризуется собственным набором атрибутов, метаданных и семантических связей, что будет детально описано в разделе 2 при проектировании логической структуры базы знаний.

\subsubsection*{Основные понятия, объекты и отношения}

Ключевыми \textit{объектами} предметной области являются:

\begin{itemize}
    \item \textbf{Концепции (CONCEPT)} — абстрактные или конкретные сущности, описывающие фундаментальные идеи Computer Science (например, \textit{«Двоичное дерево поиска»}, \textit{«Семафор»}, \textit{«Принцип инкапсуляции»}, \textit{«Протокол TCP»}). Каждая концепция характеризуется определением, областью применения, сложностными оценками и связями с другими концепциями.
    
    \item \textbf{Алгоритмы и структуры данных} — процедурные знания, описывающие способы организации и обработки информации (сортировки, обходы графов, хеширование, динамическое программирование). Для каждого алгоритма фиксируются: вход/выход, асимптотическая сложность, инварианты, типовые варианты реализации.
    
    \item \textbf{Паттерны решения задач (PATTERN)} — обобщённые стратегии применения знаний в контексте интервью (например, \textit{«Two Pointers»}, \textit{«Sliding Window»}, \textit{«DFS/BFS template»}). Паттерны связывают концепции с типовыми формулировками задач и приёмами их декомпозиции.
    
    \item \textbf{Задачи (PROBLEM)} — алгоритмические задачи с условиями, решениями на различных языках программирования, тестовыми примерами и подробными разборами. Задачи классифицируются по домену, уровню сложности (Junior/Middle/Senior) и типу.
    
    \item \textbf{Вопросы (QUESTION)} — единицы проверки теоретических знаний, включающие формулировку вопроса, эталонный ответ, критерии оценки и ссылки на релевантные концепции.
    
    \item \textbf{Шпаргалки (CHEATSHEET)} — компактные справочные материалы, содержащие ключевые формулы, определения, шаблоны кода и мнемонические правила для быстрого повторения перед собеседованием.
    
    \item \textbf{Типичные ошибки (PITFALL)} — описания распространённых заблуждений и ошибок кандидатов (например, \textit{«Путаница между стеклом и кучей»}, \textit{«Некорректная оценка сложности вложенных циклов»}) с пояснениями, как их избежать.
    
    \item \textbf{Компетенции} — абстрактные метки, отражающие уровень владения тем или иным разделом (например, \textit{«Понимает асимптотический анализ»}, \textit{«Умеет реализовывать обход дерева»}). Компетенции используются для диагностики и формирования учебных траекторий.
\end{itemize}

\textit{Отношения} между объектами формируют семантическую сеть предметной области и включают:

\begin{itemize}
    \item \textbf{Таксономические связи} (\texttt{is-a}, \texttt{part-of}): иерархия концепций (например, \textit{«Красно-чёрное дерево»} \texttt{is-a} \textit{«Самобалансирующееся дерево»});
    \item \textbf{Пререквизиты} (\texttt{requires}): зависимость понимания одной концепции от другой (например, для изучения \textit{«Динамического программирования»} необходимо знать \textit{«Рекурсию»} и \textit{«Мемоизацию»});
    \item \textbf{Применение} (\texttt{used-in}): связь алгоритма/паттерна с типами задач или доменами (например, \textit{«BFS»} \texttt{used-in} \textit{«Задачи на поиск кратчайшего пути»});
    \item \textbf{Оценка сложности} (\texttt{has-complexity}): привязка асимптотических характеристик к алгоритмам и структурам данных;
    \item \textbf{Уровень освоения} (\texttt{target-level}): соответствие концепции или задачи определённому уровню кандидата (Junior/Middle/Senior);
    \item \textbf{Проверка компетенции} (\texttt{assesses}): связь вопроса/задачи с компетенциями, которые она проверяет.
\end{itemize}

\subsubsection*{Ограничения и особенности предметной области}

При разработке базы знаний необходимо учитывать следующие \textit{ограничения и особенности}:

\begin{itemize}
    \item \textbf{Объём и динамика знаний}: предметная область характеризуется большим количеством концепций (сотни единиц) и умеренной скоростью обновления (появление новых паттернов, изменение требований компаний). База знаний должна поддерживать инкрементальное пополнение без нарушения целостности.
    
    \item \textbf{Требования к точности}: знания в области Computer Science требуют высокой формальной корректности (особенно в части асимптотического анализа, инвариантов алгоритмов, спецификаций протоколов). Недопустимы противоречивые определения или некорректные оценки сложности.
    
    \item \textbf{Требования к времени ответа}: в интерактивных сценариях (поиск, диагностика, генерация объяснений) время формирования ответа не должно превышать 1–2 секунд для сохранения качества пользовательского опыта.
    
    \item \textbf{Источники знаний}: основными источниками являются академические учебники, документация технических компаний, материалы открытых образовательных платформ, а также экспертные оценки технических интервьюеров. Все знания должны проходить верификацию и сопровождаться ссылками на авторитетные источники.
    
    \item \textbf{Мультиязычность и вариативность формулировок}: одни и те же концепции могут называться по-разному в различных источниках (например, \textit{«hash table»} / \textit{«hash map»} / \textit{«ассоциативный массив»}). База знаний должна поддерживать синонимию и нормализацию терминов.
    
    \item \textbf{Контекстная зависимость}: сложность и глубина изложения материала должны адаптироваться под роль пользователя (кандидат/ментор) и целевую компанию (стартап/FAANG).
\end{itemize}

Таким образом, предметная область «Фундаментальные основы Computer Science» представляет собой структурируемое, но объёмное множество взаимосвязанных концепций, требующее применения формальных методов представления знаний для обеспечения точности, согласованности и эффективности использования в интеллектуальной системе подготовки к собеседованиям.

\subsection{Типы задач и сценарии использования}
\label{subsec:tasks-scenarios}

Разрабатываемая база знаний фундаментальных основ Computer Science предназначена для поддержки широкого спектра интеллектуальных задач в контексте подготовки к техническим собеседованиям. Ниже приведена систематизация типов задач, решаемых с использованием БЗ, а также детализированные сценарии использования (use-cases), демонстрирующие взаимодействие пользователей и внешних модулей системы с базой знаний.

\subsubsection*{Классификация типов задач}

На основе анализа предметной области и требований целевой аудитории выделены следующие основные типы задач, для решения которых предназначена база знаний:

\begin{enumerate}
    \item \textbf{Диагностика уровня знаний} — определение степени владения пользователем ключевыми концепциями и компетенциями на основе анализа ответов на вопросы и решения задач.
    
    \item \textbf{Классификация материалов} — автоматическая категоризация единиц знаний по доменам (алгоритмы, структуры данных, ОС, сети, БД), уровням сложности (Junior/Middle/Senior) и типам (теория, задача, паттерн).
    
    \item \textbf{Поиск и выбор релевантных знаний} — семантический поиск концепций, задач и объяснений по запросам пользователя с учётом контекста (роль, уровень, целевая компания).
    
    \item \textbf{Планирование учебной траектории} — формирование персонализированной последовательности изучения тем на основе выявленных пробелов в знаниях и целевых требований позиции.
    
    \item \textbf{Объяснение концепций} — предоставление определений, примеров, визуализаций и связей с другими концепциями с варьированием глубины изложения.
    
    \item \textbf{Генерация вопросов и задач} — создание типовых формулировок интервью-вопросов и задач на код на основе шаблонов, хранящихся в БЗ.
\end{enumerate}

Каждый тип задач опирается на определённые виды знаний и механизмы доступа к ним, что будет детально описано в разделе 2 при проектировании формализма представления знаний.

\subsubsection*{Сценарии использования (Use-Cases)}

Для демонстрации практического применения базы знаний сформулированы четыре реалистичных сценария использования, охватывающих основные роли пользователей системы и типы взаимодействий.

\paragraph{Сценарий UC-1: Диагностика уровня подготовки кандидата}

\begin{itemize}
    \item \textbf{Актор:} Кандидат (Junior/Middle/Senior уровень).
    \item \textbf{Входные данные:} Профиль пользователя (целевая роль, желаемая компания), ответы на серию диагностических вопросов.
    \item \textbf{Обращение к БЗ:} Модуль диагностики запрашивает из БЗ вопросы по указанным доменам, соответствующие уровню кандидата, а также эталонные ответы и критерии оценки.
    \item \textbf{Обработка:} Система сопоставляет ответы пользователя с эталонами, вычисляет процент правильных ответов по каждой теме, определяет пробелы в знаниях.
    \item \textbf{Выходные данные:} Профиль компетенций с оценкой по каждому домену, список рекомендованных тем для изучения, визуализация прогресса.
    \item \textbf{Связанные типы задач:} Диагностика, классификация, планирование.
\end{itemize}

\paragraph{Сценарий UC-2: Персонализированная рекомендация материалов для подготовки}

\begin{itemize}
    \item \textbf{Актор:} Кандидат, осуществляющий подготовку к собеседованию.
    \item \textbf{Входные данные:} Результаты диагностики, целевая позиция (например, \textit{«Backend Developer, Middle, Yandex»}), доступное время на подготовку.
    \item \textbf{Обращение к БЗ:} Модуль рекомендаций запрашивает концепции и задачи, связанные с выявленными пробелами, фильтрует их по уровню сложности и домену, учитывает пререквизиты.
    \item \textbf{Обработка:} Система формирует упорядоченную последовательность изучения с учётом зависимостей между темами (например, перед \textit{«Динамическим программированием»} рекомендуется изучить \textit{«Рекурсию»} и \textit{«Мемоизацию»}).
    \item \textbf{Выходные данные:} Учебный план на 2–4 недели с ежедневными задачами, ссылками на материалы и оценками времени на изучение.
    \item \textbf{Связанные типы задач:} Планирование, поиск, классификация.
\end{itemize}

\paragraph{Сценарий UC-3: Объяснение концепции с адаптивной глубиной}

\begin{itemize}
    \item \textbf{Актор:} Кандидат или ментор в процессе обучения.
    \item \textbf{Входные данные:} Запрос на объяснение концепции (например, \textit{«Объясни, как работает хеш-таблица»}), указание желаемого уровня детализации (базовый/продвинутый).
    \item \textbf{Обращение к БЗ:} Модуль объяснений извлекает из БЗ карточку концепции с определением, примерами реализации, оценками сложности, связанными паттернами и типовыми вопросами.
    \item \textbf{Обработка:} Система адаптирует формулировки под уровень пользователя, выбирает релевантные примеры, формирует связный текст с перекрёстными ссылками.
    \item \textbf{Выходные данные:} Структурированное объяснение с разделами: определение, принцип работы, пример кода, сложность, связанные концепции, типовые вопросы на собеседовании.
    \item \textbf{Связанные типы задач:} Объяснение, поиск, классификация.
\end{itemize}

\paragraph{Сценарий UC-4: Генерация интервью-сессии для ментора}

\begin{itemize}
    \item \textbf{Актор:} Технический интервьюер или ментор.
    \item \textbf{Входные данные:} Параметры сессии (уровень кандидата, длительность, фокус-домены, компания-ориентир).
    \item \textbf{Обращение к БЗ:} Модуль генерации запрашивает из БЗ набор вопросов и задач, соответствующих указанным критериям, включая критерии оценки и подсказки для интервьюера.
    \item \textbf{Обработка:} Система комбинирует вопросы разных типов (теория, задача на код, системный дизайн), обеспечивает покрытие ключевых тем, балансирует сложность \cite{Xu2020}.
    \item \textbf{Выходные данные:} Готовый сценарий интервью на 45–60 минут с таймингом, вопросами, ожидаемыми ответами и рекомендациями по оценке.
    \item \textbf{Связанные типы задач:} Генерация, классификация, поиск.
\end{itemize}

\subsubsection*{Связь сценариев с архитектурой системы}

Описанные сценарии использования формируют не только общее представление о поведении системы, но и задают конкретные требования к проектированию интерфейсов доступа к базе знаний, включая структуру входных данных, формат возвращаемых результатов и правила взаимодействия между модулями. На основании этих сценариев уточняются механизмы обработки запросов: как выполняется их первичная интерпретация, каким образом определяется тип задачи, какие компоненты участвуют в маршрутизации запроса и как обеспечивается согласованность ответов при различных пользовательских контекстах. Кроме того, сценарии позволяют выявить критические точки, в которых необходимы дополнительные проверки корректности, устойчивости и производительности, чтобы система сохраняла предсказуемое поведение при изменении характера нагрузки. В таблице~\ref{tab:use-cases-summary} приведено обобщённое соответствие между выделенными сценариями, типами решаемых задач и задействованными компонентами системы, что позволяет наглядно сопоставить функциональные требования с архитектурной реализацией и показать, каким образом отдельные элементы системы совместно обеспечивают выполнение целевых пользовательских процессов.

\begin{table}[h]
\centering
\caption{Соответствие сценариев использования, типов задач и компонентов системы}
\label{tab:use-cases-summary}
\small
\setlength{\tabcolsep}{4pt}
\begin{tabularx}{\linewidth}{|>{\raggedright\arraybackslash}p{0.23\linewidth}|>{\raggedright\arraybackslash}p{0.27\linewidth}|>{\raggedright\arraybackslash}X|}
\hline
\textbf{Сценарий} & \textbf{Типы задач} & \textbf{Задействованные компоненты} \\ \hline
UC-1: Диагностика & Диагностика, классификация & Модуль диагностики, БЗ вопросов \\ \hline
UC-2: Рекомендации & Планирование, поиск & Модуль рекомендаций, БЗ концепций \\ \hline
UC-3: Объяснение & Объяснение, поиск & Модуль объяснений, БЗ концепций \\ \hline
UC-4: Генерация интервью & Генерация, классификация & Модуль генерации, БЗ задач \\ \hline
\end{tabularx}
\end{table}

\subsubsection*{Ожидаемые результаты использования}

Реализация описанных сценариев с использованием разрабатываемой базы знаний должна обеспечить следующие результаты:

\begin{itemize}
    \item \textbf{Для кандидатов:} сокращение времени подготовки на 30–40\% за счёт персонализации, повышение уверенности за счёт структурированного покрытия тем, объективная оценка готовности к собеседованию.
    
    \item \textbf{Для менторов и интервьюеров:} стандартизация процесса проведения технических интервью, снижение нагрузки на подготовку материалов, повышение согласованности оценки кандидатов.
    
    \item \textbf{Для системы в целом:} масштабируемая основа для расширения предметных областей (System Design, Soft Skills), возможность интеграции с внешними образовательными платформами, поддержка аналитики эффективности подготовки.
\end{itemize}

Таким образом, сценарии использования демонстрируют практическую ценность базы знаний как центрального компонента интеллектуальной системы подготовки, обеспечивающего решение ключевых задач диагностики, рекомендаций, объяснения и генерации материалов для технических собеседований.

\subsection{Анализ существующих решений и баз знаний}
\label{subsec:existing-solutions}

Для обоснования архитектурных и формальных решений при разработке базы знаний «Фундаментальные основы Computer Science» был проведён сравнительный анализ существующих платформ и систем, решающих сходные задачи представления знаний и поддержки подготовки к техническим собеседованиям.

\subsubsection*{Обзор аналогичных систем}

В качестве аналогов выбраны три системы, представляющие различные парадигмы организации знаний: платформа оценки навыков, академическая онтологическая система и визуально-ориентированный образовательный ресурс.

\paragraph{HackerRank}

\textbf{Назначение:} Платформа для оценки технических компетенций и проведения скрининга кандидатов, широко используемая технологическими компаниями на ранних этапах отбора \cite{web:hackerrank}.

\textbf{Тип базы знаний:} Гибридная документно-ориентированная база с элементами графа навыков (\textit{skill graph}).

\textbf{Формы представления знаний:}
\begin{itemize}
    \item JSON-структуры для описания задач, тестовых кейсов, решений и метаданных;
    \item Иерархическая модель компетенций: \textit{язык программирования → домен → тема → подтема → задача};
    \item Связи «задача → требуемый навык → минимальный уровень» для адаптивного подбора контента.
\end{itemize}

\paragraph{Computer Science Ontology (CSO Portal)}

\textbf{Назначение:} Академический проект по созданию формальной онтологии области компьютерных наук для семантической аннотации, поиска и анализа научных публикаций и образовательных ресурсов \cite{web:csoportal}.

\textbf{Тип базы знаний:} Онтологическая база знаний на основе стандартов W3C (RDF/OWL) с поддержкой логического вывода.

\textbf{Формы представления знаний:}
\begin{itemize}
    \item Формальные классы и свойства (\texttt{cso:Algorithm}, \texttt{cso:hasTimeComplexity}, \texttt{cso:requiresPrerequisite});
    \item Иерархия концепций с использованием \texttt{rdfs:subClassOf}, \texttt{owl:equivalentClass}, \texttt{owl:disjointWith};
    \item Связи с внешними источниками (DBpedia, Wikidata) через \texttt{owl:sameAs} для обогащения контекста.
\end{itemize}

\paragraph{VisuAlgo}

\textbf{Назначение:} Интерактивная образовательная платформа для визуализации алгоритмов и структур данных, ориентированная на улучшение понимания через анимацию и практическое исполнение \cite{web:visualgo}.

\textbf{Тип базы знаний:} Концептуальный граф знаний (\textit{concept graph}) с мультимодальными компонентами (текст, код, визуализация).

\textbf{Формы представления знаний:}
\begin{itemize}
    \item Узлы графа представляют концепции (\textit{«Binary Heap»}, \textit{«Dijkstra's Algorithm»});
    \item Рёбра кодируют отношения: \texttt{requires} (пререквизиты), \texttt{visualizes} (связь с анимацией), \texttt{implemented-in} (язык реализации);
    \item Привязка псевдокода, реализаций на популярных языках и пошаговых анимаций к узлам графа.
\end{itemize}

\subsubsection*{Анализ структур баз знаний существующих решений}

В таблице~\ref{tab:kb-alternatives-comparison} приведено сравнительное описание структур баз знаний рассмотренных систем по ключевым критериям.

\begin{table}[h]
\centering
\caption{Сравнительный анализ структур баз знаний аналогов}
\label{tab:kb-alternatives-comparison}
\small
\setlength{\tabcolsep}{4pt}
\begin{tabularx}{\linewidth}{|>{\raggedright\arraybackslash}p{0.22\linewidth}|>{\raggedright\arraybackslash}X|>{\raggedright\arraybackslash}X|>{\raggedright\arraybackslash}X|}
\hline
\textbf{Критерий} & \textbf{HackerRank} & \textbf{CSO Portal} & \textbf{VisuAlgo} \\ \hline
\textbf{Виды знаний} & Задачи, навыки, оценки & Концепции, отношения, аксиомы & Алгоритмы, визуализации, код \\ \hline
\textbf{Логическая организация} & Skill graph + теги + метаданные & Формальная онтология OWL 2 DL & Концептуальный граф + мультимедиа \\ \hline
\textbf{Физическая организация} & Реляционная БД + JSON-документы & RDF-триплы + SPARQL-endpoint & Файловая система + JS-граф \\ \hline
\textbf{Механизмы вывода} & Статистические рекомендации, фильтрация по профилю & Дедуктивный вывод через OWL-reasoner & Навигация по связям, интерактивное исполнение \\ \hline
\textbf{Полнота} & Высокая по задачам, средняя по теории & Высокая по академическим терминам & Высокая по визуализированным алгоритмам \\ \hline
\textbf{Согласованность} & Обеспечивается модерацией, без формальной верификации & Формальная верификация через логические ограничения & Частичная, зависит от ручной синхронизации \\ \hline
\textbf{Расширяемость} & Ограничена внутренней схемой платформы & Высокая (стандарты W3C, линковка) & Средняя (требует ручной доработки визуализаций) \\ \hline
\textbf{Сопровожда-
емость} & Централизованная команда, закрытый процесс & Сообщество исследователей, открытый репозиторий & Небольшая команда разработчиков \\ \hline
\end{tabularx}
\end{table}

\newpage

\textbf{Сильные и слабые стороны рассмотренных решений:}

\begin{itemize}
    \item \textbf{HackerRank:}
    \begin{itemize}
        \item \textit{Сильные стороны:} богатая коллекция практических задач, интеграция с системами найма, адаптивная подборка на основе профиля навыков;
        \item \textit{Слабые стороны:} слабая формализация теоретических знаний, отсутствие явных семантических связей между концепциями, закрытость структуры БЗ для внешнего использования.
    \end{itemize}
    
    \item \textbf{CSO Portal:}
    \begin{itemize}
        \item \textit{Сильные стороны:} строгая формальная модель, поддержка логического вывода и проверки непротиворечивости, совместимость с Semantic Web-стандартами;
        \item \textit{Слабые стороны:} ориентация на академические публикации, а не на подготовку к собеседованиям, высокий порог входа для авторов контента.
    \end{itemize}
    
    \item \textbf{VisuAlgo:}
    \begin{itemize}
        \item \textit{Сильные стороны:} уникальная интеграция декларативных знаний с визуальными и процедурными компонентами, высокая педагогическая эффективность;
        \item \textit{Слабые стороны:} ограниченное покрытие тем (только визуализируемые алгоритмы), отсутствие механизмов адаптации под уровень пользователя.
    \end{itemize}
\end{itemize}

\textbf{Выводы для проектирования базы знаний проекта:}

На основе проведённого анализа принято решение использовать \textit{комбинированный подход}, заимствующий лучшие практики рассмотренных систем:

\begin{enumerate}
    \item В качестве логической основы принять элементы онтологической модели, аналогичной CSO Portal (бинарные отношения, таксономия, ограничения целостности), для обеспечения формальной согласованности и поддержки семантического поиска;
    
    \item Для представления контента использовать человекочитаемый формат Markdown с YAML-фронтматтером, что упростит наполнение и сопровождение по аналогии с практиками HackerRank;
    
    \item Реализовать механизмы визуальной навигации и привязки примеров кода к концепциям, вдохновлённые подходом VisuAlgo, для повышения педагогической эффективности материалов.
\end{enumerate}

Такой гибридный подход позволит совместить формальную строгость академических онтологий с практической удобностью коммерческих платформ и наглядностью образовательных ресурсов, что полностью соответствует целям разрабатываемого интеллектуального помощника по подготовке к техническим собеседованиям.

\subsection{Анализ требований к базе знаний}

\subsubsection*{Функциональные требования}

% Требования:
% - сформулировать, какие задачи должна обеспечивать база знаний: какие запросы и выводы она должна поддерживать;
% - связать каждое требование с типом задач и сценариями использования;
% - использовать формулировки «база знаний должна обеспечивать …».

На основании анализа типов задач и сценариев использования сформулированы следующие функциональные требования для предметной области «Фундаментальные основы Computer Science»:

\begin{enumerate}
    \item База знаний должна обеспечивать диагностику уровня владения алгоритмическими компетенциями путём сопоставления ответов пользователя с карточками задач (\texttt{PROBLEM}) и вопросами (\texttt{QUESTION}) с использованием связей типа \texttt{evaluates} и матриц компетенций (\texttt{MATRIX}).
    \item База знаний должна обеспечивать построение индивидуальных учебных траекторий на основе обхода графа знаний, где рёбра \texttt{prerequisite\_for} и \texttt{depends\_on} определяют допустимую последовательность изучения концепций (например, \texttt{Массив} $\rightarrow$ \texttt{Динамическое программирование}).
    \item База знаний должна обеспечивать семантический поиск и фильтрацию единиц знаний по метаданным (домен \texttt{cs}, категория \texttt{concept}/\texttt{problem}/\texttt{pattern}, уровень сложности, роль, этап интервью, теги) с поддержкой AND-логики и ранжирования по релевантности.
    \item База знаний должна обеспечивать предоставление объяснений концепций с адаптивной глубиной изложения через извлечение данных из карточек типов \texttt{CONCEPT} и \texttt{CHEATSHEET} с учётом уровня пользователя (\texttt{junior}/\texttt{middle}/\texttt{senior}).
    \item База знаний должна обеспечивать кросс-доменную интеграцию знаний путём поддержки связей между предметными областями (\texttt{cs/}, \texttt{arch/}, \texttt{process/}) и формирования единого реестра идентификаторов (\texttt{GLB-CROSS}), например, связь алгоритма \texttt{Consistent Hashing} с задачей системного дизайна \texttt{URL Shortener}.
\end{enumerate}\par

\subsubsection*{Нефункциональные требования}

% Требования:
% - сформулировать требования к полноте, согласованности, точности, расширяемости и сопровождаемости базы знаний;
% - указать ограничения по объёмам данных, производительности обработки, средствам разработки и хранения;
% - использовать лексику «следует», «необходимо», «не допускается».

Для обеспечения качества и долгосрочной жизнеспособности системы необходимо соблюдение следующих требований:

\begin{enumerate}
    \item \textit{Полнота покрытия.} Объём базы знаний по ПО «Фундаментальные основы CS» должен составлять не менее 200 единиц знаний, включая концепции, задачи, паттерны, вопросы, шпаргалки и типичные ошибки, согласно установленным минимумам ТЗ.
    \item \textit{Согласованность и непротиворечивость.} Терминология должна строго соответствовать глобальному глоссарию. Дублирование идентификаторов, циклические зависимости в графе предпосылок и противоречивые оценки сложности (\texttt{Big O}) в связанных документах не допускаются.
    \item \textit{Точность и актуальность.} Каждая единица знаний должна подтверждаться минимум одним авторитетным источником. Асимптотические оценки, инварианты алгоритмов и спецификации протоколов должны быть формально корректны; дата последнего обновления (\texttt{date\_updated}) для материалов по трендам не должна превышать 2 года. Подход к фиксации инженерной точности и проверяемости знаний учитывает практики архитектурного анализа и описания надёжных систем \cite{Kleppmann2021}. 
    \item \textit{Расширяемость и сопровождаемость.} Структура хранения должна позволять добавление новых единиц знаний без модификации ядра системы. Использование унифицированного формата (Markdown + YAML-фронтматтер) необходимо для обеспечения версионирования, автоматизированной валидации и последующей интеграции с решателями задач.
    \item \textit{Производительность обработки.} Операции поиска, фильтрации и валидации перекрёстных ссылок должны выполняться в интерактивном режиме (до 2 секунд для выборки до 50 единиц) при использовании штатного CLI-инструмента.
\end{enumerate}\par

\subsubsection*{Требования к видам знаний и формализму}

% Требования:
% - описать требуемые виды знаний: декларативные (факты, таксономии, отношения), процедурные (правила вывода, алгоритмические схемы), ограничения целостности;
% - зафиксировать требования к формализму представления знаний (логика, онтологии, продукционные правила, форматы данных и т.п.);
% - указать ожидания по выразительной мощности, удобству обработки и интеграции с решателем задач и интерфейсами.

База знаний должна поддерживать репрезентативный набор видов знаний, релевантных задачам подготовки к техническим собеседованиям в области Computer Science:
\begin{itemize}
    \item \textit{Декларативные знания:} факты о свойствах структур данных (\texttt{CONCEPT}), таксономии алгоритмических парадигм (\texttt{PATTERN}), описания вычислительной сложности и ограничений целостности (инварианты, предусловия/постусловия).
    \item \textit{Процедурные знания:} алгоритмические схемы решения типовых задач, паттерны декомпозиции (\texttt{Sliding Window}, \texttt{Two Pointers}), чек-листы анализа сложности (\texttt{CHEATSHEET}), критерии оценки корректности кода. Подход к выделению повторно используемых шаблонов решения согласуется с практиками паттерн-ориентированного проектирования \cite{Gamma2020}. 
    \item \textit{Метазнания и ограничения:} правила валидации структуры документов, матрицы соответствия навыков уровням сложности, учебные траектории (\texttt{TRAJECTORY}), правила калибровки сложности задач (\texttt{LeetCode Easy/Medium/Hard} $\leftrightarrow$ \texttt{junior/middle/senior}).
\end{itemize}\par

В качестве основного формализма представления знаний выбран полуструктурированный документно-ориентированный подход на базе Markdown с YAML-фронтматтером. Данный выбор обоснован следующими факторами:
\begin{enumerate}
    \item \textit{Выразительная мощность:} YAML обеспечивает строгую типизацию метаданных (идентификаторы формата \texttt{CS-CATEGORY-NN}, домены, теги, связи), а Markdown позволяет хранить развёрнутые объяснения, ASCII-диаграммы структур данных и примеры кода на псевдокоде/Python, что критично для алгоритмической предметной области.
    \item \textit{Удобство обработки:} Формат парсится стандартными библиотеками (\texttt{pyyaml}), что позволяет реализовать автоматизированную валидацию, генерацию реестров и семантический поиск без развёртывания тяжёлых СУБД.
    \item \textit{Интеграционный потенциал:} Документная структура совместима с конвейерами RAG (векторизация, чанкинг), легко конвертируется в RDF/OWL при необходимости перехода к онтологическому формализму, и полностью поддерживается системами контроля версий (Git) для совместной работы трёх инженеров знаний.
\end{enumerate}\par
Не допускается смешение несовместимых формализмов без явной схемы согласования. Все связи между документами должны быть явно зафиксированы в полях \texttt{related} и \texttt{prerequisites} с указанием типа связи в комментариях (например, \texttt{\# depends\_on}, \texttt{\# applied\_in}).

\subsection{Вывод по разделу анализа}

% Требования:
% - в 3–6 пунктах обобщить результаты анализа предметной области, типов задач, аналогов и требований;
% - явно указать, какие требования и решения будут использованы при проектировании базы знаний;
% - не вводить новой информации.

На основании проведённого исследования можно сформулировать следующие итоги, определяющие направление проектирования базы знаний для предметной области «Фундаментальные основы Computer Science»:

\begin{enumerate}
    \item Предметная область требует представления знаний, ориентированного не только на декларативные описания алгоритмов, но и на процедурные паттерны решения задач (\texttt{PATTERN}), анализ типичных ошибок (\texttt{PITFALL}) и адаптацию глубины изложения под уровень кандидата.
    \item Существующие образовательные платформы демонстрируют ограничения в формализации семантических связей между концепциями, поддержке персонализированной диагностики алгоритмических навыков и глубокой проработке компромиссов (time vs. space complexity), что обосновывает необходимость разработки специализированной базы знаний.
    \item База знаний должна решать задачи диагностики компетенций, планирования учебных траекторий, семантического поиска и генерации объяснений с варьированием детализации, что требует внедрения графовой модели отношений (\texttt{prerequisite\_for}, \texttt{depends\_on}, \texttt{applied\_in}, \texttt{warns\_about}).
    \item Функциональные требования диктуют необходимость поддержки фильтрации по метаданным (домен, категория, сложность, роль, этап интервью), обхода графа зависимостей и кросс-доменной интеграции с предметными областями ARCH и PROCESS.
    \item Нефункциональные требования фиксируют минимальный объём ($\geq$200 единиц для ПО «CS»), строгую согласованность терминологии, обязательную верификацию источников и использование полуструктурированного формата Markdown+YAML для обеспечения машиночитаемости, расширяемости и совместной работы.
    \item В проектировочном разделе будет обоснован выбор документно-графовой архитектуры БЗ, реализована схема валидации метаданных через CLI-инструмент, определены форматы входных/выходных данных для модулей диагностики и рекомендаций, и разработаны диаграммы компонентов системы, соответствующие сформулированным требованиям.
\end{enumerate}